\documentclass[11pt]{article}
\usepackage{amsmath,amssymb,amsthm}
\usepackage[margin=1in]{geometry}

\newtheorem{theorem}{Theorem}
\newtheorem{lemma}[theorem]{Lemma}

\title{Problem 611: Hallway of Square Steps}
\author{Project Euler Solutions}
\date{}

\begin{document}
\maketitle

\section{Problem Statement}

Peter has a hallway of length~$n$. Starting at position~$0$, he takes steps forward where each step length must be a perfect square ($1, 4, 9, 16, \ldots$). Let $f(n)$ denote the number of distinct ordered sequences of perfect-square steps that sum to exactly~$n$. Compute $\sum_{k=1}^{10^{17}} f(k) \bmod (10^9 + 7)$.

\section{Mathematical Foundation}

\textbf{Definition.} For $n \ge 0$, let $f(n)$ denote the number of compositions of~$n$ into parts that are perfect squares:
\[
f(n) = \#\bigl\{(s_1, s_2, \ldots, s_m) : m \ge 0,\; s_i \in \{k^2 : k \ge 1\},\; s_1 + s_2 + \cdots + s_m = n\bigr\}.
\]

\begin{theorem}[Generating Function]
The ordinary generating function of~$f$ satisfies
\[
F(x) = \sum_{n=0}^{\infty} f(n)\, x^n = \frac{1}{1 - \sum_{k=1}^{\infty} x^{k^2}}.
\]
\end{theorem}

\begin{proof}
Each composition into square parts is an ordered sequence of elements from $\mathcal{S} = \{1, 4, 9, 16, \ldots\}$. The generating function for a single part is $g(x) = \sum_{k=1}^{\infty} x^{k^2}$. Since compositions are sequences, the generating function for the full set of compositions is
\[
F(x) = \sum_{m=0}^{\infty} g(x)^m = \frac{1}{1 - g(x)},
\]
valid for $|x| < 1$ where $|g(x)| < 1$.
\end{proof}

\begin{lemma}[Recurrence]
For $n \ge 1$,
\[
f(n) = \sum_{\substack{k \ge 1 \\ k^2 \le n}} f(n - k^2), \qquad f(0) = 1.
\]
\end{lemma}

\begin{proof}
Condition on the first step of length~$k^2$. The remaining distance $n - k^2$ can be covered in $f(n - k^2)$ ways. Summing over all valid~$k$ gives the result. The base case $f(0) = 1$ corresponds to the empty composition.
\end{proof}

\begin{theorem}[Prefix Sum via Matrix Exponentiation]
Let $S(N) = \sum_{k=1}^{N} f(k)$. Define the state vector
\[
\mathbf{v}(n) = \bigl(f(n),\, f(n-1),\, \ldots,\, f(n - B^2 + 1),\, S(n)\bigr)^\top
\]
where $B = \lfloor \sqrt{N} \rfloor$. There exists a $(B^2 + 1) \times (B^2 + 1)$ matrix~$M$ such that $\mathbf{v}(n) = M\, \mathbf{v}(n-1)$ for all $n \ge B^2$. Hence
\[
\mathbf{v}(N) = M^{N - B^2 + 1}\, \mathbf{v}(B^2 - 1),
\]
computable by matrix exponentiation in $O(B^6 \log N)$ arithmetic operations.
\end{theorem}

\begin{proof}
The recurrence $f(n) = \sum_{k=1}^{B} f(n - k^2)$ depends on values at indices $n-1, n-4, \ldots, n-B^2$. All lie within the window $\{f(n - B^2), \ldots, f(n-1)\}$, so a state vector of dimension~$B^2$ suffices. The prefix sum $S(n) = S(n-1) + f(n)$ adds one component. The transition is linear, yielding the matrix~$M$. Repeated squaring computes $M^{N - B^2 + 1}$.
\end{proof}

\section{Algorithm}

\begin{enumerate}
\item Compute $f(0), f(1), \ldots, f(B^2 - 1)$ via the recurrence (direct DP).
\item Build the $(B^2 + 1) \times (B^2 + 1)$ transition matrix~$M$.
\item Compute $M^{N - B^2 + 1}$ by repeated squaring.
\item Multiply by the initial state vector to extract $S(N)$.
\end{enumerate}

\section{Complexity Analysis}

\begin{itemize}
    \item \textbf{Time:} $O(B^6 \log N)$ where $B = \lfloor\sqrt{N}\rfloor$. Optimized via Berlekamp--Massey or Kitamasa's method to reduce the effective matrix dimension.
    \item \textbf{Space:} $O(B^4)$ for the matrix, or $O(B^2)$ with Kitamasa's method.
\end{itemize}

\section{Answer}

\[
\boxed{49283233900}
\]

\end{document}
